\section{Introduction}
\label{sec:intro}

Deep learning models fine-tuned on diverse downstream tasks often exhibit complementary capabilities. Composing these task-specific experts into a single unified multi-task model is a fundamental challenge in modern machine learning engineering. The conventional approach of joint multi-task learning (MTL) is often computationally prohibitive and logistically impractical, requiring simultaneous access to all private downstream datasets, massive training compute, and costly hyperparameter optimization \cite{caruana1997multitask, standley2020tasks}.

To circumvent the cost of joint training, \emph{parameter-space model merging} has emerged as a powerful alternative. Pioneers like Task Arithmetic \cite{ilharco2023editing} and Model Soups \cite{wortsman2022model} demonstrate that we can simply subtract the pre-trained base model's weights from the fine-tuned experts' weights to form task vectors, and then linearly average these vectors back into the base weights. Because this weight-space composition happens post-hoc, it requires no access to training datasets or joint backpropagation. Crucially, the resulting merged model retains the exact architecture and parameter size of the base model, introducing \textbf{exactly zero inference latency or parameter overhead} during edge deployment.

Despite these advantages, simple linear merging methods suffer from severe performance degradation due to \emph{inter-task interference}. Because weight updates in different fine-tuned models can conflict, simply averaging them can damage critical functional representations. To address this, state-of-the-art adaptive model merging frameworks like AdaMerging \cite{yu2024adamerging}, SyMerge \cite{lu2026symerge}, and FoldMerge have been developed. These frameworks optimize layer-wise or parameter-wise merging coefficients at test-time, guided by test-time adaptation (TTA) losses or predictions from teacher networks.

\begin{figure}[t]
  \centering
  \includegraphics[width=0.48\textwidth]{cost_accuracy_tradeoff.png}
  \caption{\textbf{The Merging Pareto Frontier (Preparation Time vs.\ 8-Task Average Accuracy)} on CLIP ViT-B/32. Under a pragmatic lens, state-of-the-art adaptive methods like SyMerge \cite{lu2026symerge} require joint gradient-based optimization over 500 steps, taking over 10 minutes on high-end GPUs and requiring enormous gradient memory. EdgeMerge (Ours) completely bypasses backpropagation, extracting adaptive channel-wise coefficients in just 11.95 seconds on a single forward pass, matching the peak optimized Task Arithmetic baseline (68.74\%) while requiring near-zero GPU memory.}
  \label{fig:pareto}
\end{figure}

However, through a \emph{pragmatic engineering lens}, these state-of-the-art adaptive methods are deeply impractical for real-world production deployment. They are characterized by several major limitations:
\begin{itemize}
    \item \textbf{High Adaptation Latency:} Running 500 gradient steps takes up to 10+ minutes even on server-grade hardware (such as an NVIDIA H100 GPU). For an edge device or resource-constrained environment (e.g., embedded systems, IoT nodes, or CPUs), waiting 10 minutes before the first multi-task prediction can be made is completely unacceptable.
    \item \textbf{Severe Memory Footprint:} Gradient-tracking during test-time backpropagation requires maintaining activation graphs and multiple copies of large teacher/student models in memory. Most edge hardware has severe memory limits that immediately trigger out-of-memory (OOM) faults.
    \item \textbf{The Overfitting-Optimizer Paradox:} Recent analyses reveal that optimizing numerous fine-grained layer coefficients on small test batches is highly prone to transductive overfitting, where delicate test-time tuning on a narrow subset of data degrades the model's generalizability on broader, unseen data.
\end{itemize}

To bridge the gap between adaptive performance and deployment practicality, we propose \textbf{EdgeMerge}, a completely training-free, forward-only adaptive model merging framework. EdgeMerge is designed specifically to operate within real-world engineering constraints: it completely removes backpropagation, runs in seconds, and uses minimal memory, making it highly robust and deployable on edge devices. Importantly, we resolve a common logical inconsistency in edge-calibration literature: while our algorithm is mathematically lightweight enough to run within tight on-device resource budgets (e.g., by sequentially streaming expert checkpoints from local flash memory to perform the forward pass), its primary and most practical real-world engineering workflow is \emph{offline calibration}. A developer can run our 11.95-second calibration pass on a local workstation using small validation samples, reconstruct the single merged multi-task checkpoint, and ship it to edge hardware. This workflow completely bypasses the need for on-device checkpoint storage or test-time calibration while ensuring that the deployed model is optimized and conflict-free. 

Our core methodology revolves around computing localized, channel-wise merging coefficients in closed-form from a single forward pass of a tiny, unlabeled calibration dataset (e.g., 32 samples per task). EdgeMerge consists of three main stages: (1) \textbf{Forward-Only Activation Sampling (FOAS)} to capture the internal functional behaviors of the experts and base model; (2) \textbf{Scale-Normalized Delta Activation Salience (SNDAS)} to estimate the importance of each neuron/channel for each task without bias from natural activation scales; and (3) \textbf{Channel-Wise Softmax Gating (CWSG)} to normalize importance across tasks, yielding fine-grained, localized merging coefficients that resolve inter-task interference.

We evaluate EdgeMerge on the official 8-task visual classification benchmark using the modular Vision-Language CLIP ViT-B/32 architecture. Following the rigorous standards of our pragmatic philosophy, we do not compare against a weak baseline; instead, we thoroughly optimize the Task Arithmetic (TA) baseline by performing a grid search over global scaling factors $\lambda \in [0.1, 0.8]$ to establish a strong, empirical lower bound.

Our experimental findings highlight a highly compelling trade-off:
\begin{enumerate}
    \item \textbf{Extreme Efficiency Gains:} EdgeMerge completely eliminates backpropagation and gradient-tracking. It computes adaptive channel-wise merging weights in only \textbf{11.95 seconds}---representing an over \textbf{$50\times$ speedup} compared to SyMerge (600 seconds)---while keeping training GPU memory overhead at approximately 100 MB.
    \item \textbf{Transparent Performance-Accuracy Trade-off:} Our forward-only, training-free approach incurs a substantial accuracy penalty of \textbf{21.05\% absolute points} compared to server-grade, fully optimized gradient-based methods like SyMerge (68.69\% vs. 89.74\%). This frames EdgeMerge as an extreme-efficiency exploration for resource-constrained edge systems rather than a raw accuracy competitor under unconstrained server-side conditions.
    \item \textbf{Hyperparameter Robustness \& Stabilization:} Standard Task Arithmetic is highly fragile, showing a sharp, narrow peak in performance at $\lambda = 0.20$ and collapsing rapidly elsewhere. In contrast, by dynamically routing visual projection bottleneck channels, EdgeMerge accommodates a much more robust and stable scaling range (e.g., maintaining high performance at $\lambda = 0.30$), mitigating the major deployment risks of static weight averaging.
\end{enumerate}

By showing that a single forward pass can achieve the same peak performance as an optimized global grid search while stabilizing the scaling hyperparameter space, EdgeMerge demonstrates that training-free, forward-only activation routing is a highly viable, practical path forward for scalable and robust model merging on resource-constrained edge hardware.
